\section{Object Model}

MIFX represents machining data as a collection of independent objects stored as JSON files within the package.

Each object describes a specific aspect of the machining process and may reference additional data through artifacts.

This design allows systems to interpret only the parts of the package they understand while safely ignoring unsupported objects or extensions.

\subsection{Primary Objects}

The core MIFX object types are:

\begin{itemize}
\item \textbf{Job} — the root object describing the machining package
\item \textbf{Setup} — a manufacturing setup configuration
\item \textbf{Operation} — an individual machining operation
\item \textbf{Tool} — a cutting tool definition
\item \textbf{In-Process} — an optional intermediate workpiece state
\end{itemize}

Each object is stored as a JSON file and may reference additional resources through artifacts.

\subsection{Object Identity}

Each object must define an identifier using the \texttt{id} field.

Example:

\begin{verbatim}
{
  "id": "op-2"
}
\end{verbatim}

Object identity is defined by the \texttt{id} field rather than by the filename.  
File names within the package are therefore not required to follow a specific naming convention.

\subsection{Object References}

Objects may reference other objects or resources using relative file paths.

Example:

\begin{verbatim}
{
  "id": "op-2",
  "artifacts": [
    {
      "role": "tool",
      "kind": "json",
      "path": "../../tools/main/2.json",
      "present": true
    }
  ]
}
\end{verbatim}

This approach allows the package to behave as a lightweight object graph where nodes reference related resources.

\subsection{Object Independence}

Objects within a MIFX package are designed to be loosely coupled.

Software implementations may interpret only the objects and artifacts relevant to their functionality while ignoring unsupported elements.

This design enables interoperability between systems with different capabilities while maintaining forward compatibility.